\documentclass[aps,prd,twocolumn,superscriptaddress,nofootinbib]{revtex4-2}

\usepackage[T1]{fontenc}
\usepackage{amsmath,amssymb}
\usepackage{booktabs}
\usepackage{microtype}
\usepackage{siunitx}
\usepackage{xcolor}
\usepackage{hyperref}
\usepackage{pgfplots}
\pgfplotsset{compat=1.18}
\sisetup{range-phrase=--,range-units=single}
\setcounter{dbltopnumber}{3}
\hypersetup{
  colorlinks=true,
  linkcolor=blue!45!black,
  citecolor=green!35!black,
  urlcolor=blue!55!black
}

\definecolor{tracedarkblue}{HTML}{005A8D}
\definecolor{tracelightblue}{HTML}{56B4E9}
\definecolor{adjointdarkorange}{HTML}{C45100}
\definecolor{adjointlightorange}{HTML}{E69F00}
\definecolor{multipletgreen}{HTML}{009E73}

\begin{document}

\title{Exact colour sums for multi-gluon amplitudes:\texorpdfstring{\\}{ }
direct, multiplet and Fourier methods}

\author{Rikkert Frederix}
\email{rikkert.frederix@fysik.lu.se}
\affiliation{Department of Physics, Lund University, Box 118, 221 00 Lund, Sweden}
\author{Valentin Hirschi}
\email{hirschi@itp.unibe.ch}
\affiliation{Institute for Theoretical Physics, University of Bern,
Sidlerstrasse 5, 3012 Bern, Switzerland}
\author{Malin Sj\"odahl}
\email{malin.sjodahl@fysik.lu.se}
\affiliation{Department of Physics, Lund University, Box 118, 221 00 Lund, Sweden}

\begin{abstract}
We compare exact $SU(3)$ colour sums for fixed-helicity, tree-level all-gluon
matrix elements.  In the trace and adjoint decompositions, a direct colour
contraction couples every pair of partial amplitudes, so its evaluation time
grows quadratically with the factorially growing number of gluon orderings.  An
orthonormal multiplet basis removes this final double sum but requires a
representation-labelled amplitude recursion.  Alternatively, a fast Fourier
transform (FFT) on the symmetric group retains the trace or adjoint partial
amplitudes and reorganizes the exact colour sum into smaller contractions
labelled by Young diagrams.
We compare CPU time after initialization and maximum memory use for four
through eleven gluons.  At eleven gluons, one matrix-element evaluation takes
\SI{19.8}{s} in the multiplet basis, \SI{3.08}{s} with the trace FFT, and
\SI{0.814}{s} with the adjoint FFT; the direct trace and adjoint calculations
are estimated to take \SI{2.41e5}{s} and \SI{1.69e3}{s}.  The adjoint FFT is
the fastest method from six through eleven gluons.
\end{abstract}

\maketitle

\section{Introduction}
\label{sec:intro}

Colour-ordered Berends--Giele recursion efficiently computes tree-level gluon
partial amplitudes~\cite{Berends:1987me}.  Computing an exact colour-summed
squared matrix element also requires evaluating all interference terms among
the associated colour structures.
At fixed external momenta and helicities, let $\mathcal M$ denote the full
amplitude.  The index $\alpha$ labels the chosen colour tensors $C_\alpha$,
and $A_\alpha$ is the corresponding complex kinematic
coefficient:
\begin{equation}
  \mathcal M=\sum_{\alpha} C_{\alpha}\,A_{\alpha}.
  \label{eq:colour-decomposition}
\end{equation}
The sum over all external colours is then
\begin{equation}
  \begin{aligned}
    \mathcal B
    &\equiv\sum_{\text{colours}}|\mathcal M|^2
    =\sum_{\alpha,\beta}A_{\alpha}^{*}G_{\alpha\beta}A_{\beta},\\
    G_{\alpha\beta}
    &\equiv\sum_{\text{colours}}C_{\alpha}^{*}C_{\beta}.
  \end{aligned}
  \label{eq:colour-sum}
\end{equation}
The star denotes complex conjugation, the vertical bars denote the complex
modulus, and $\alpha$ and $\beta$ both label the colour tensors in
Eq.~\eqref{eq:colour-decomposition}.  The quantity $\mathcal B$ is the
colour-summed squared matrix element, while $G$ is the colour-overlap matrix
with entries $G_{\alpha\beta}$.  The structure of $G$, and hence the
work required to evaluate Eq.~\eqref{eq:colour-sum}, depends on the chosen
colour decomposition.

The trace decomposition~\cite{Paton:1969je,Mangano:1987xk} and the adjoint
basis of Del Duca, Dixon and Maltoni~\cite{DelDuca:1999rs} express the
amplitude in terms of familiar colour-ordered partial amplitudes.  Their
colour tensors are not orthogonal, so the off-diagonal terms in
Eq.~\eqref{eq:colour-sum} must be retained.  In both cases, the number of
orderings grows factorially with the gluon multiplicity, making direct
evaluation of the double sum increasingly expensive.

Multiplet bases offer a different organization of colour
space~\cite{Keppeler:2012ih,Sjodahl:2018cca}.  They are constructed by coupling
gluons through irreducible $SU(3)$ representations, and their colour tensors
can be chosen orthonormal.  The final colour sum then reduces to a sum of
absolute squares.  This simplification requires the recursion to track
representation labels and recouple the corresponding colour states when
gluons are reordered.  Here we present the first numerical implementation of
the off-shell multiplet recursion
developed in Ref.~\cite{Bolinder:2025gbj}.  We also use this recursion for the
first time to evaluate fully colour-summed squared matrix elements at fixed
momenta and helicities.

The non-orthogonal trace and adjoint decompositions nevertheless have useful
structure.  In either case, the overlap between two colour tensors depends
only on the relative permutation of their gluon orderings.  A fast Fourier
transform (FFT) on the symmetric group exploits this property to replace the
direct double sum with a set of smaller, independent matrix contractions.
The Fourier method is exact: it retains every colour interference term and
does not change the partial amplitudes.  To our knowledge, this is the first
application of a symmetric-group FFT to exact colour summation for QCD matrix
elements.

We examine these three approaches in turn.  Section~\ref{sec:direct}
describes direct colour summation in the trace and adjoint decompositions.
Section~\ref{sec:multiplet} presents the orthogonal multiplet basis and its
recursion.  Section~\ref{sec:fft} explains the symmetric-group Fourier method
for the trace and adjoint colour sums.  Section~\ref{sec:results} compares
their CPU times per evaluation after initialization and their maximum memory
use.  Throughout, $N$ denotes the total number of external gluons.

\section{Direct colour sums in trace and adjoint decompositions}
\label{sec:direct}

\subsection{Trace decomposition}

Cyclic invariance lets us fix one external gluon in the single-trace
decomposition.  For any positive integer $q$, $S_q$ denotes the group of
permutations of $q$ labels; $\sigma\in S_q$ is one such permutation, and
$\sigma(i)$ is the image of label $i$.  Permuting the remaining $N-1$ gluons
gives
\begin{equation}
  \mathcal M
  =\sum_{\sigma\in S_{N-1}}
   C^{\mathrm{tr}}_{\sigma}
   A\bigl(\sigma(1),\ldots,\sigma(N-1),N\bigr),
  \label{eq:trace-decomposition}
\end{equation}
with
\begin{equation}
  P_{\mathrm{tr}}=(N-1)!
  \label{eq:trace-size}
\end{equation}
partial amplitudes, where $!$ denotes factorial.  Here
$C^{\mathrm{tr}}_\sigma$ is the trace colour tensor for order $\sigma$,
$A(\sigma(1),\ldots,\sigma(N-1),N)$ is the corresponding partial amplitude,
and $P_{\mathrm{tr}}$ is the number of trace orders.  The trace set is
overcomplete but convenient because its kinematic coefficients are the usual
colour-ordered amplitudes.  Reflection relates each ordering to its reverse,
so the corresponding currents can be shared when evaluating partial
amplitudes.

\subsection{Adjoint basis}

The adjoint decomposition places gluons $1$ and $N$ at the ends of a chain of
structure constants and permutes the remaining gluons~\cite{DelDuca:1999rs}:
\begin{equation}
  \mathcal M
  =\sum_{\sigma\in S_{N-2}}
   C^{\mathrm{adj}}_{\sigma}
   A\bigl(1,\sigma(2),\ldots,\sigma(N-1),N\bigr).
  \label{eq:adjoint-decomposition}
\end{equation}
There are
\begin{equation}
  P_{\mathrm{adj}}=(N-2)!
  \label{eq:adjoint-size}
\end{equation}
partial amplitudes.  Here $C^{\mathrm{adj}}_\sigma$ is the chain of structure
constants for order $\sigma$,
$A(1,\sigma(2),\ldots,\sigma(N-1),N)$ is the corresponding partial amplitude,
and $P_{\mathrm{adj}}$ is the number of adjoint orders.  The adjoint set is
smaller than the trace set, but its colour tensors remain non-orthogonal.

\subsection{Direct contraction}

For either decomposition, let $P$ denote $P_{\mathrm{tr}}$ or
$P_{\mathrm{adj}}$, as appropriate.  For trace, let $\sigma$, $\tau$, and
$\pi$ belong to $S_{N-1}$; for adjoint, let them belong to $S_{N-2}$, which
here permutes labels $2,\ldots,N-1$.  Direct evaluation of
Eq.~\eqref{eq:colour-sum} contracts one pair of orders at a time.  Because the
colour matrix is real and symmetric, we evaluate each diagonal term once,
while each off-diagonal pair contributes twice the real part of its
upper-triangular term.  A common relabelling of both colour tensors leaves
their overlap unchanged, so
\begin{equation}
  G_{\sigma\tau}=k(\sigma^{-1}\tau),
  \qquad k(\pi)=k(\pi^{-1}).
  \label{eq:relative-permutation}
\end{equation}
Here $G_{\sigma\tau}$ is the overlap of the colour tensors for the two orders,
the superscript $-1$ denotes permutation inversion, and $k(\pi)$ is the
overlap obtained by placing the first tensor in the reference order and using
$\pi$ as the relative permutation.  Only the $P$ values of $k$ need to be
stored because every row of $G$ is a permutation of this list.  The
contraction still considers $P^2$ pairs, so its time scales as $[(N-1)!]^2$
for trace and $[(N-2)!]^2$ for adjoint.  These direct contractions provide the
baseline for the comparison below.

For both decompositions, we use Berends--Giele
recursion~\cite{Berends:1987me} to evaluate the ordered amplitudes.  We compute
currents shared by several complete orders once and reuse them.  Our comparison
therefore includes both partial-amplitude evaluation and final colour
contraction, rather than the contraction alone.

\section{The multiplet basis}
\label{sec:multiplet}

Each gluon transforms in the adjoint representation of $SU(3)$.  To construct
a multiplet basis, we couple the first two gluons to an irreducible
representation and then couple each remaining gluon in turn until the total
representation is a singlet.  Each basis vector is labelled by the resulting
representation sequence and an extra label whenever a representation occurs
more than once.  With a suitable normalization, the colour tensors obey
\begin{equation}
  \sum_{\text{colours}}
  C_{\alpha}^{*}C_{\beta}=\delta_{\alpha\beta}.
  \label{eq:multiplet-orthogonality}
\end{equation}
Here $\alpha$ and $\beta$ label basis vectors, including any extra labels for
repeated representations, and $\delta_{\alpha\beta}$ is the Kronecker delta.
Equation~\eqref{eq:colour-sum} therefore becomes
\begin{equation}
  \mathcal B=\sum_{\alpha}|A_{\alpha}|^2.
  \label{eq:multiplet-sum}
\end{equation}
Thus, once the amplitude coefficients $A_\alpha$ are known, exact colour
summation reduces to an ordinary sum of absolute squares.

We compute the amplitude coefficients in this basis using the off-shell
recursion of Ref.~\cite{Bolinder:2025gbj}.  A current is labelled by a subset
of external gluons and by an allowed sequence of $SU(3)$ representations.
Joining two currents requires reordering their gluons into the order defining
the parent basis.  We decompose this reordering into exchanges of neighbouring
gluons.  Each exchange acts on the representation sequences through a
normalized matrix of Wigner $6j$ coefficients~\cite{Sjodahl:2015qoa}.

These coefficients depend on the normalization, phase, and multiplicity
conventions chosen for the three-point $SU(3)$ vertices.  We calculate these
coefficients in the multiplet-basis conventions and tabulate the resulting
normalized exchange matrices.  We validate the complete table in four ways: by
direct contractions of explicit adjoint tensors, by contractions after
resolving each adjoint into a traceless fundamental--antifundamental pair, and
by recursive reduction to the elementary two-line recouplings described in
Refs.~\cite{Alcock-Zeilinger:2022hrk,Keppeler:2023msu}, and by a
carrier-space-free construction using analytic fusion labels, dimensions,
Casimirs, and Young-seminormal swaps.  Direct symmetric-group projector
calculations provide a further independent check of the seed and representative
multiplicity-free and repeated-channel blocks.  The resulting
tabulation contains 354 exchange blocks with 5673 nonzero entries and covers
every recoupling required at the multiplicities considered.  Because the
coefficients and exchange sequences depend only on $N$, we prepare them once
before evaluating phase-space points.

The representations allowed in repeated adjoint products determine the
number of multiplet coefficients.  For $N=4,\ldots,11$, the respective counts
are $8$, $32$, $145$, $702$, $3598$, $19280$, $107160$, and $614000$.
The multiplet basis is smaller than the trace set from $N=7$ onward but larger
than the adjoint basis throughout this range.
Orthogonality eliminates the final double sum, whereas the
representation-labelled recursion requires more work than the ordered trace
and adjoint recursions.  The timings below show whether eliminating the double
sum offsets this additional work.

\section{Symmetric-group FFT for trace and adjoint colour sums}
\label{sec:fft}

For either decomposition, let $P=m!$, where $m=N-1$ for trace and $m=N-2$
for adjoint.  The partial amplitudes define a function $A(\sigma)$ on the
permutation group $S_m$.  By Eq.~\eqref{eq:relative-permutation}, each colour
overlap depends only on the relative permutation.  The colour sum is therefore
a convolution on $S_m$.  Specifically, for each $\tau\in S_m$, the
convolution sums $A(\sigma)k(\sigma^{-1}\tau)$ over $\sigma\in S_m$; the
colour sum then contracts the result with $A(\tau)^*$ and sums over $\tau$.

On a cyclic group, a Fourier transform converts convolution into scalar
multiplication.  Because permutations do not generally commute, the
corresponding transform on $S_m$ instead produces matrices.  Its blocks are
labelled by partitions $\lambda$ of $m$, or equivalently by Young diagrams
with $m$ boxes.  We write $\lambda\vdash m$ for such a partition.  If
$d_\lambda$ is the dimension of the
irreducible representation labelled by $\lambda$, then
\begin{equation}
  \sum_{\lambda\vdash m}d_\lambda^2=m!.
  \label{eq:young-counting}
\end{equation}
Thus the $m!$ amplitude values are reorganized exactly into one
$d_\lambda\times d_\lambda$ matrix for each Young diagram.

For a complex-valued function $f$ on $S_m$, define its transformed matrix
$\widehat f_\lambda$ by
\begin{equation}
  \widehat f_\lambda
  =\sum_{\sigma\in S_m}f(\sigma)\rho_\lambda(\sigma),
  \label{eq:group-transform}
\end{equation}
where $\rho_\lambda(\sigma)$ is the real orthogonal
$d_\lambda\times d_\lambda$ matrix representing the permutation $\sigma$ in
the irreducible representation labelled by $\lambda$.  Let
$\widehat A_\lambda$ be the transform of the partial-amplitude function
$A(\sigma)$, and let $K_\lambda$ be the transform of the colour-overlap
function $k(\sigma)$.  The exact colour sum is
\begin{equation}
  \mathcal B=\frac{1}{m!}
  \sum_{\lambda\vdash m}d_\lambda
  \operatorname{Tr}\!\left(
    \widehat A_\lambda^\dagger
    \widehat A_\lambda K_\lambda
  \right).
  \label{eq:fft-colour-sum}
\end{equation}
The dagger denotes Hermitian conjugation.
Equation~\eqref{eq:fft-colour-sum} couples only matrices associated with the
same Young diagram, rather than all pairs of permutations.

Computing every transformed matrix entry directly from
Eq.~\eqref{eq:group-transform} would still be expensive.  Fast transforms on
finite groups instead reuse partial results~\cite{DiaconisRockmore1990}.  For
the symmetric group, we evaluate the transform recursively along the subgroup
chain
\begin{equation}
  S_1\subset S_2\subset\cdots\subset S_m,
  \label{eq:subgroup-chain}
\end{equation}
where $S_i$ is the permutation group on $i$ labels for $i=1,\ldots,m$.  For
$i=1,\ldots,m-1$, the inclusion $S_i\subset S_{i+1}$ leaves label $i+1$
unchanged.  Partial results at one level are reused at the next.  Along this
chain, the matrices representing exchanges of adjacent gluons split into
one- and two-dimensional blocks, reducing the work at each
step~\cite{ClausenBaum1993,Maslen1998}.

For each $N$, we calculate the $K_\lambda$ matrices once and reuse them.
Collectively, these matrices contain
$P=\sum_{\lambda\vdash m}d_\lambda^2$ colour entries.  At each new phase-space
point and helicity assignment, we transform the partial amplitudes and
evaluate the contractions in Eq.~\eqref{eq:fft-colour-sum}.  The transform
does not make either set of colour tensors orthogonal or reduce the number of
partial amplitudes; it reorganizes the exact colour sum to avoid contracting
$P^2$ pairs directly.

In our trace calculation, we obtain the colour overlaps by closing fundamental
colour lines in $U(3)$.  Although individual $U(3)$ overlaps need not equal
their $SU(3)$ counterparts, the $U(1)$ component decouples from a tree-level
pure-gluon amplitude.  Contracting these $U(3)$ overlaps with the physical
partial amplitudes therefore gives the same final $SU(3)$ colour sum.  For the
adjoint decomposition, we write the chains of structure constants as nested
commutators.  Their $U(1)$ components vanish because the identity generator
commutes with every generator.  The trace and adjoint cases use different
values of $m$ and different overlap functions $k$, but the same
symmetric-group transform.

\section{Numerical comparison}
\label{sec:results}

\begin{table*}[t]
  \caption{CPU time for one evaluation of a fixed-helicity, colour-summed
  squared matrix element after initialization.  Initialization and the first
  evaluation are excluded.  Direct denotes colour-overlap-matrix contraction
  without a Fourier transform, in either the trace decomposition or the
  adjoint basis.  FFT denotes the exact symmetric-group Fourier method.
  Asterisks mark estimates, not measurements.}
  \label{tab:evaluation-times}
  \centering
  \setlength{\tabcolsep}{5.5pt}
  \begin{tabular}{@{}c rr r rr@{}}
  \toprule
  $N$ & \multicolumn{2}{c}{direct colour sum} &
  \multicolumn{1}{c}{orthogonal basis} &
  \multicolumn{2}{c}{Fourier colour sum} \\
  \cmidrule(lr){2-3}\cmidrule(lr){4-4}\cmidrule(lr){5-6}
  & trace & adjoint & multiplet & trace FFT & adjoint FFT \\
  \midrule
  4  & \SI{0.729}{\micro\second} & \SI{0.543}{\micro\second} & \SI{0.923}{\micro\second} & \SI{0.968}{\micro\second} & \SI{0.645}{\micro\second} \\
  5  & \SI{6.23}{\micro\second}  & \SI{2.44}{\micro\second}  & \SI{6.38}{\micro\second}  & \SI{4.41}{\micro\second}  & \SI{2.62}{\micro\second} \\
  6  & \SI{138}{\micro\second}   & \SI{15.3}{\micro\second}  & \SI{62.3}{\micro\second}  & \SI{26.2}{\micro\second}  & \SI{13.3}{\micro\second} \\
  7  & \SI{6.16}{\milli\second}  & \SI{0.196}{\milli\second} & \SI{0.718}{\milli\second} & \SI{0.186}{\milli\second} & \SI{0.0832}{\milli\second} \\
  8  & \SI{0.381}{\second}       & \SI{5.78}{\milli\second}  & \SI{8.26}{\milli\second}  & \SI{1.62}{\milli\second}  & \SI{0.620}{\milli\second} \\
  9  & \SI{29.7}{\second}        & \SI{0.303}{\second}       & \SI{0.104}{\second}       & \SI{0.0155}{\second}      & \SI{0.00523}{\second} \\
  10 & \SI{2.41e3}{\second}\textsuperscript{*} & \SI{20.9}{\second} & \SI{1.49}{\second} & \SI{0.215}{\second} & \SI{0.0589}{\second} \\
  11 & \SI{2.41e5}{\second}\textsuperscript{*} & \SI{1.69e3}{\second}\textsuperscript{*} & \SI{19.8}{\second} & \SI{3.08}{\second} & \SI{0.814}{\second} \\
  \bottomrule
  \end{tabular}

  \vspace{0.5ex}
  \begin{minipage}{0.97\textwidth}
  \footnotesize
  \textsuperscript{*}Estimated from the last measured direct result, assuming
  that the time is proportional to the square of the number of
  colour orders, with $P_{\mathrm{tr}}=(N-1)!$ and
  $P_{\mathrm{adj}}=(N-2)!$.
  \end{minipage}
\end{table*}


\subsection{Setup and checks}

We study $gg\to(N-2)g$, summing the colours of all incoming and outgoing
gluons while holding their helicities fixed.

For each $N=4,\ldots,11$, we apply all five methods to the same phase-space
point and the same helicity assignment, chosen to give a nonzero matrix
element.  We use the general Berends--Giele recursion in every case, even when
an analytic expression is available, so special-case formulas do not affect
the comparison.  At fixed $N$, the recursion performs the same operations for
every helicity assignment not excluded by the tree-level selection rule.

We compiled the five implementations with GNU Fortran 13.3 at optimization
level \texttt{-O3} and ran each on a single core of an Intel Core i7-8700K.
Timing excludes initialization and the first evaluation, during which any
remaining $N$-dependent data are prepared.  The reported CPU time per
subsequent matrix-element evaluation is the median of ten measurements.

For each method, we measured maximum physical memory in a separate process;
the values include initialization and all evaluations.

For every tested configuration, all completed calculations agree on the
colour-summed squared matrix element.  The largest pairwise relative
difference between nonzero results is $8.1\times10^{-12}$, at ten gluons.  For
helicity configurations forbidden by the tree-level selection rule, all five
methods return exactly zero.  Extrapolated times are excluded from these
checks.

\subsection{Evaluation time after initialization}

Table~\ref{tab:evaluation-times} and Fig.~\ref{fig:evaluation-times} summarize
the timing results.  Direct adjoint contraction is fastest at four and five
gluons; adjoint with the Fourier transform is fastest from six through eleven
gluons.  The direct double sum becomes costly by eight gluons: direct trace
takes \SI{0.381}{s}, whereas trace with the transform takes \SI{1.62}{ms}.
At nine gluons, the corresponding times are \SI{29.7}{s} and \SI{15.5}{ms},
a factor of about $1900$.

The multiplet method avoids the direct double sum.  It is faster than direct
trace at six gluons and faster than direct adjoint at nine gluons.  At eleven
gluons, the multiplet evaluation takes \SI{19.8}{s}, versus an estimated
\SI{1.69e3}{s} for direct adjoint and measured times of \SI{3.08}{s} and
\SI{0.814}{s} for trace and adjoint with the transform, respectively.  The
adjoint Fourier method is therefore fastest at this multiplicity.

\begin{figure}[t]
  \centering
  \begin{tikzpicture}
    \begin{axis}[
      width=0.97\columnwidth,
      height=0.70\columnwidth,
      xlabel={number of external gluons $N$},
      ylabel={CPU time per evaluation [s]},
      xmin=3.7,xmax=11.3,
      xtick={4,5,6,7,8,9,10,11},
      ymode=log,
      ymin=1e-7,ymax=1e6,
      ymajorgrids=true,
      grid style={draw=black!12,line width=0.35pt},
      axis line style={draw=black!65,line width=0.5pt},
      tick align=outside,
      legend style={at={(0.02,0.98)},anchor=north west,
                    draw=black!18,fill=white,fill opacity=0.94,
                    text opacity=1,font=\scriptsize,rounded corners=1pt,
                    legend columns=2,
                    /tikz/every even column/.append style={column sep=7pt}},
      tick label style={font=\scriptsize},
      label style={font=\small},
      unbounded coords=jump,
    ]
      \addplot+[mark=square*,mark size=2.3pt,thick,densely dashed,tracedarkblue,
                mark options={solid,fill=tracedarkblue,draw=tracedarkblue}]
        table[x=N,y=trace_direct] {data/colour_sum_times.dat};
      \addlegendentry{trace, direct}
      \addplot+[mark=diamond*,mark size=2.5pt,thick,tracelightblue,
                mark options={solid,fill=tracelightblue,draw=tracelightblue}]
        table[x=N,y=trace_fft] {data/colour_sum_times.dat};
      \addlegendentry{trace, FFT}
      \addplot+[mark=triangle*,mark size=2.6pt,thick,densely dashed,adjointdarkorange,
                mark options={solid,fill=adjointdarkorange,draw=adjointdarkorange}]
        table[x=N,y=adjoint_direct] {data/colour_sum_times.dat};
      \addlegendentry{adjoint, direct}
      \addplot+[mark=pentagon*,mark size=2.4pt,thick,adjointlightorange,
                mark options={solid,fill=adjointlightorange,draw=adjointlightorange}]
        table[x=N,y=adjoint_fft] {data/colour_sum_times.dat};
      \addlegendentry{adjoint, FFT}
      \addplot+[mark=*,mark size=2.2pt,thick,multipletgreen,
                mark options={solid,fill=multipletgreen,draw=multipletgreen}]
        table[x=N,y=multiplet] {data/colour_sum_times.dat};
      \addlegendentry{multiplet}
      \addplot+[no marks,thick,densely dotted,tracedarkblue,forget plot]
        table[x=N,y=trace_extension] {data/colour_sum_times.dat};
      \addplot+[no marks,thick,densely dotted,adjointdarkorange,forget plot]
        table[x=N,y=adjoint_extension] {data/colour_sum_times.dat};
      \addplot+[only marks,mark=square,mark size=2.7pt,
                mark options={fill=white,draw=tracedarkblue,line width=0.8pt},forget plot]
        table[x=N,y=trace_estimate] {data/colour_sum_times.dat};
      \addplot+[only marks,mark=triangle,mark size=3.0pt,
                mark options={fill=white,draw=adjointdarkorange,line width=0.8pt},forget plot]
        table[x=N,y=adjoint_estimate] {data/colour_sum_times.dat};
    \end{axis}
  \end{tikzpicture}
  \caption{CPU time to evaluate one fixed-helicity, colour-summed squared
  matrix element after initialization.  Dashed lines show direct trace and
  adjoint colour sums; solid blue and orange lines show their exact Fourier
  counterparts; the solid green line shows the multiplet method.  Dotted
  extensions and open symbols indicate the three estimates identified by
  asterisks in Table~\ref{tab:evaluation-times}, while filled symbols indicate
  measurements.}
  \label{fig:evaluation-times}
\end{figure}

\subsection{Memory use}

The following table gives maximum memory use at the four largest
multiplicities.  A dash indicates that the corresponding run did not finish;
no memory value is estimated.
\begin{center}
  \small
  \setlength{\tabcolsep}{3.2pt}
  \begin{tabular}{@{}l rrrr@{}}
  \toprule
  method & $N=8$ & $N=9$ & $N=10$ & $N=11$ \\
  \midrule
  trace, direct & \SI{4.24}{MiB} & \SI{14.0}{MiB} & --- & --- \\
  adjoint, direct & \SI{3.34}{MiB} & \SI{7.66}{MiB} & \SI{43.9}{MiB} & --- \\
  multiplet     & \SI{10.1}{MiB} & \SI{55.4}{MiB} & \SI{425}{MiB} & \SI{3.96}{GiB} \\
  trace, FFT    & \SI{4.43}{MiB} & \SI{14.8}{MiB} & \SI{111}{MiB} & \SI{1.09}{GiB} \\
  adjoint, FFT  & \SI{3.54}{MiB} & \SI{7.86}{MiB} & \SI{44.6}{MiB} & \SI{398}{MiB} \\
  \bottomrule
  \end{tabular}
\end{center}


At every multiplicity with a direct result, the direct and Fourier variants
use similar amounts of memory.  At nine gluons, direct trace uses
\SI{14.0}{MiB} and trace with the transform uses
\SI{14.8}{MiB}; the adjoint values are \SI{7.66}{MiB} and \SI{7.86}{MiB}.
Our direct implementations store only the $P$ distinct overlaps from
Eq.~\eqref{eq:relative-permutation}, not the full $P\times P$ matrix.  Within
each decomposition, the direct and Fourier methods therefore differ far more
in evaluation time than in stored colour data.

At eleven gluons, memory use is \SI{3.96}{GiB} for the multiplet method,
\SI{1.09}{GiB} for trace with the transform, and \SI{398}{MiB} for adjoint
with the transform.  For each of the three methods with results at both
multiplicities, memory use grows by about an order of magnitude from ten to
eleven gluons, suggesting that memory may become an important limitation
beyond the range studied here.

\section{Conclusions}
\label{sec:conclusion}

We have compared three exact approaches to evaluating fixed-helicity,
colour-summed squared matrix elements for tree-level all-gluon scattering.
Although our direct trace and adjoint implementations each store only one row
of colour overlaps, both contractions couple every pair of partial amplitudes
and therefore scale quadratically with the factorially growing number of
amplitudes.

The multiplet-basis calculation reduces the final colour sum to a sum of
absolute squares and substantially outperforms the direct contractions at
high multiplicity, but the symmetric-group Fourier method is faster still.
It retains the trace or adjoint partial amplitudes while exploiting the
relative-permutation dependence of their colour overlaps to replace the
direct double sum with exact contractions labelled by Young diagrams.

At eleven gluons, the CPU times per evaluation after initialization are
\SI{3.08}{s} for trace with the transform and \SI{0.814}{s} for adjoint with
the transform, compared with \SI{19.8}{s} in the multiplet basis.  Both
Fourier methods also use less memory at this multiplicity, despite the
multiplet basis's simpler final colour sum.  The direct trace and adjoint
contractions are estimated to require \SI{2.41e5}{s} and \SI{1.69e3}{s},
respectively.  The adjoint Fourier method uses the smaller $(N-2)!$ set
of partial amplitudes; it is fastest from six through eleven gluons and has
the smallest recorded memory use at eleven gluons.

Processes with external quarks are a natural next application.  The
pure-gluon adjoint chain cannot be carried over unchanged because quarks and
antiquarks carry fundamental and antifundamental colour.  For one
quark--antiquark pair, the corresponding colour decomposition uses an open
chain of fundamental generators, with the gluons permuted between the quark
and antiquark.  For several pairs, the established generalization also tracks
the placement and nesting of the quark
lines~\cite{Johansson:2015oia,Melia:2015ika}.  Multiplet bases can accommodate
external quarks, antiquarks, and gluons while remaining orthogonal, but
extending the recursion requires representation paths and recoupling
coefficients involving fundamental
lines~\cite{Sjodahl:2018cca,Alcock-Zeilinger:2022hrk,Keppeler:2023msu}.  For
one pair, the open-chain colour overlaps again depend only on the relative
gluon ordering, so the same symmetric-group transform applies once the
required open-chain overlaps have been computed.  With several pairs, one
could apply the transform to gluon permutations within each quark-line
arrangement, but different arrangements remain coupled; whether such partial
transforms reduce the total cost remains to be determined.

\bibliographystyle{apsrev4-2}
\bibliography{references}

\end{document}
